\documentclass[11pt,letterpaper]{article}
\usepackage[T1]{fontenc}
\usepackage[utf8]{inputenc}
\usepackage[letterpaper,textwidth=13.9cm,top=1.45cm,bottom=1.6cm,footskip=25pt]{geometry}
\usepackage{amsmath,booktabs,tabularx,array}
\usepackage[hidelinks]{hyperref}
\setlength{\parskip}{4pt}
\setlength{\intextsep}{9pt}
\linespread{1.0}
\setlength{\parindent}{0pt}
\setlength{\emergencystretch}{1.8em}
\renewcommand{\arraystretch}{1.13}
\newcolumntype{Y}{>{\raggedright\arraybackslash}X}
\newcommand{\code}[1]{\nolinkurl{#1}}
\newcommand{\fonte}[1]{\par\smallskip\noindent{\fontsize{8}{10}\selectfont\textit{Source.} #1\par}}
\newcommand{\pagina}[2]{\section{#2}}
\makeatletter
\renewcommand\section{\@startsection{section}{1}{\z@}{2.5ex plus .4ex}{1.3ex}{\normalfont\normalsize\bfseries}}
\renewcommand\subsection{\@startsection{subsection}{2}{\z@}{1.8ex plus .3ex}{.8ex}{\normalfont\normalsize\itshape}}
\def\ps@primeira{\def\@oddhead{}\def\@oddfoot{\small\itshape Empirical study\hfil November 2026}}
\makeatother
\hypersetup{pdftitle={Volatility Forecasting and Position Sizing in BOVA11}}

\begin{document}
\thispagestyle{primeira}
\vspace*{.65cm}
\begin{center}
{\Large\bfseries Volatility Forecasting and Position Sizing\\[3pt] in BOVA11\par}
\vspace{1.7em}
{\itshape Daily evidence and execution costs\par}
\end{center}
\vspace{1.15em}\hrule\smallskip
\noindent\textbf{Abstract}\par
\begingroup\small
We examine whether more accurate volatility forecasts lead to better position sizing in BOVA11. Over 437 common daily targets, HAR had the lowest QLIKE in 2025 and 2026. Its mean difference against EWMA was $-0.0645$, with a 95\% block-bootstrap interval of $[-0.1457,0.0027]$ and a Holm-adjusted two-sided Diebold--Mariano $p$-value of 0.7543. In a panel of 430 intraday sessions, forecasting risk over the traded window moved annual volatility closer to the 12\% target but reduced net return from $-53.02\%$ to $-60.61\%$. On the frozen XGBoost signal, a fixed-size control at matched exposure reproduced nearly all the drawdown reduction from regime sizing. The comparisons distinguish the cost of raising exposure from the effect of lowering it on drawdown.

\textit{Keywords.} BOVA11, volatility, market regimes, QLIKE, trading costs.
\par\endgroup\smallskip\hrule\medskip

\section{Introduction}
A better volatility forecast can lead to a portfolio with worse net returns. The mechanism is straightforward when the forecast changes exposure in a strategy that opens and closes a position every day. Volatility describes the magnitude of a move. The return earned on a position also depends on direction, exposure, and execution costs.

We examine three hypotheses. $H_1$ asks whether regime models improve extreme-event forecasts against simple baselines. $H_2$ asks whether regime sizing adds value at matched exposure. $H_3$ concerns whether an economic benefit survives costs and holds prospectively. Variance forecasts provide a separate diagnostic for understanding the sizing decisions. The use of latent states follows the regime-switching literature [1].

\section{Data and methodology}
\subsection{Calendar, targets, and samples}
The origin $t$ is a BVMF trading session. For risk forecasting we use adjusted closing prices, with log return $r_{t+1}=\log(C_{t+1}^{\rm adj}/C_t^{\rm adj})$ and daily variance proxy $v_{t+1}=r_{t+1}^2$. The missing session on 24 October 2025 stays in the calendar, and returns spanning the gap are undefined. Execution uses $R^{OC}_{t+1}=C_{t+1}/O_{t+1}-1$. Execution prices for 26 October 2026 were quarantined because providers disagreed on the opening price. The adjusted close remains available for risk analysis.

\begin{table}[h!]
\centering\caption{Evaluation samples and benchmarks.}
\begin{tabularx}{\textwidth}{@{}p{3.2cm}rY@{}}
\toprule
Sample & Observations & Benchmark or control \\
\midrule
Daily variance & 437 & EWMA on common dates, 2 Jan 2025 to 5 Nov 2026 \\
Markov Switching & 438 & Fixed and refitted models compared with EWMA \\
Intraday execution & 430 & Close-to-close EWMA, fixed exposure, and a separate passive benchmark \\
Frozen signal & 10 & Matched exposure and within-side permutations on the same trades \\
\bottomrule
\end{tabularx}
\end{table}

\subsection{Models and time ordering}
The historical baseline averages 21 squared returns. EWMA uses $h_{t+1}=0.94h_t+0.06r_t^2$. HAR adapts the multi-horizon structure of Corsi [2] to a daily proxy. For $m_j(t)=j^{-1}\sum_{i=0}^{j-1}v_{t-i}$, we fit
\[
 \widehat h^{\rm HAR}_{t+1|t}
 =\max\{\varepsilon,\beta_0+\beta_1m_1(t)+\beta_5m_5(t)+\beta_{22}m_{22}(t)\}.
\]
XGBoost adds local and external features. The lab also compares GARCH, EGARCH, GJR-GARCH, and a two-state HMM. In the regime panel, $p_{t+1|t}=p_{t|t}A$, with transition matrix $A$ and filtered state probabilities $p_{t|t}$. The second-moment forecast is $\sum_k p_{t+1|t,k}(\sigma_k^2+\mu_k^2)$. The Markov switching walk-forward version refits every 63 origins when at least 126 returns are available.

Each HAR and XGBoost fit uses only targets observed by the first close of its block. S\&P 500, USD/BRL, VIX, crude oil, and EWZ are matched by dates strictly before the origin. The historical files record dates, while actual first publication times are unavailable.

\subsection{Loss and paired comparisons}
For $h>0$, QLIKE is $\ell(v,h)=\log h+v/h$ [3]. We compare $\Delta_t=\ell(v_{t+1},h^{\rm model}_{t+1|t})-\ell(v_{t+1},h^{\rm EWMA}_{t+1|t})$ on common dates. The exploratory 95\% interval uses 3,000 paired bootstrap replicates with blocks of ten sessions. Blocks preserve the boundaries of contiguous runs. The annual split uses the target date.

\section{Volatility forecasting and temporal stability}
HAR had the lowest QLIKE across the 437 common targets, followed by GARCH(1,1) and the 21-session historical baseline. In 2026, XGBoost had a larger loss relative to EWMA than it did in 2025. The ranking is therefore sensitive to the evaluation window.

\begin{table}[h!]
\centering\caption{QLIKE loss and mean differences against EWMA.}\label{tab:qlike}
\begin{tabular*}{\textwidth}{@{\extracolsep{\fill}}lrrrr@{}}
\toprule
Model & Full QLIKE & Full diff. & 2025 diff. & 2026 diff. \\
\midrule
HAR, daily proxy & $-7.685$ & $-0.064$ & $-0.047$ & $-0.084$ \\
GARCH(1,1) & $-7.636$ & $-0.016$ & $-0.033$ & $+0.002$ \\
Historical 21 & $-7.633$ & $-0.012$ & $-0.021$ & $-0.003$ \\
EWMA & $-7.621$ & $0$ & $0$ & $0$ \\
Two-state HMM & $-7.529$ & $+0.091$ & $+0.007$ & $+0.182$ \\
GJR-GARCH & $-7.500$ & $+0.120$ & $-0.016$ & $+0.266$ \\
XGBoost & $-7.486$ & $+0.135$ & $-0.045$ & $+0.328$ \\
EGARCH & $-7.368$ & $+0.253$ & $+0.002$ & $+0.522$ \\
\bottomrule
\end{tabular*}
\fonte{Local audit \code{P2.22_estabilidade_diaria.json}. The annual samples contain 226 and 211 targets.}
\end{table}

\subsection{HAR versus EWMA}
We use the Diebold--Mariano statistic [4], $DM=\overline\Delta/\widehat{\mathrm{SE}}_{\mathrm{HAC}}(\overline\Delta)$. Newey--West variance uses Bartlett weights and ten exchange-session lags. The calendar preserves the distance between observations across gaps. Holm adjustment covers seven models compared against EWMA.

For HAR, $DM=-1.608$ and the two-sided $p$-value is 0.1078, or 0.7543 after Holm. The exploratory one-sided $p$-value in favor of HAR is 0.0539, or 0.3772 after Holm. The two-sided adjusted comparisons did not establish superiority over EWMA at 5\%. HAR was selected after inspecting the historical ranking, so these tests are exploratory.

The 95\% bootstrap interval is $[-0.1436,-0.0041]$ with five-session blocks, $[-0.1457,0.0027]$ with ten, and $[-0.1543,0.0073]$ with twenty. The unadjusted two-sided HAC $p$-value ranges from 0.0437 at zero lags to 0.1229 at twenty. The inference depends on the treatment of temporal dependence.
\fonte{\code{analise/comparacao_har_ewma.json}.}

\subsection{Regimes and extreme events}
Refitted Markov switching obtained QLIKE $-7.573$, versus $-7.569$ for the fixed model and $-7.625$ for EWMA across 438 origins. Both intervals against EWMA include zero. For $H_1$, an extreme is $|r_{t+1}|>q$, with $q$ set at the training 90th percentile of $|r|$. The forecast is $\sum_k p_{t+1|t,k}P_k(|r|>q)$.

There were 35 extreme outcomes. The refitted model's Brier score was 0.07418, versus 0.07392 for the training-frequency baseline. The difference was 0.00026, with a 95\% interval from $-0.00187$ to $0.00244$. In the separate tail exercise, GARCH and the historical baseline were compared using FZ0, a joint VaR and Expected Shortfall loss. Intervals included zero at both risk levels. The incremental value of regimes for these tail measures remains an open question.
\fonte{\code{walk_forward_qlike.csv}, \code{calibracao_evento_extremo.csv}, and \code{P1_comparacao_VaR_CVaR.csv}.}

\section{Position sizing, execution costs, and the frozen signal}
\subsection{Volatility targeting}
The exposure rule is $w_t=\min\{1,\,0.12/\sqrt{252\widehat h_{t+1|t}}\}$. With signed exposure $a_t=s_tw_t$, per-side cost $c=0.0015$, and initial wealth $W_0=1$,
\[
 L_{t+1}=a_tR^{OC}_{t+1}-2c|a_t|,\qquad W_{t+1}=W_t(1+L_{t+1}).
\]
Drawdown is the decline from the running peak, including initial wealth. The annualized mean-variance certainty-equivalent return [5] is
\[
 \widehat{\mathrm{CER}}_{\mathrm{ann}}(\gamma)
 =252(\overline L-\tfrac{\gamma}{2}s_L^2).
\]
Here $\overline L$ and $s_L^2$ are the sample mean and variance of daily net returns. We report $\gamma=3$. CER uses an arithmetic annual scale, while gross and net returns are compounded over the full sample.

\begin{table}[h!]
\centering\caption{Intraday policies on 430 common forecast dates.}
\setlength{\tabcolsep}{3pt}
\small
\begin{tabular*}{\textwidth}{@{\extracolsep{\fill}}lrrrrrr@{}}
\toprule
Policy & OC QLIKE & Ann. vol. & Exposure & Gross & Net & Ann. CER \\
\midrule
Close EWMA & $-8.3252$ & $10.11\%$ & $70.82\%$ & $+17.20\%$ & $-53.02\%$ & $-45.26\%$ \\
Intraday EWMA & $-8.3930$ & $12.19\%$ & $85.14\%$ & $+18.26\%$ & $-60.61\%$ & $-56.03\%$ \\
Fixed at 100\% & --- & $14.19\%$ & $100.00\%$ & $+23.73\%$ & $-65.99\%$ & $-65.13\%$ \\
\bottomrule
\end{tabular*}
\fonte{\code{H2_previsao_e_controles.csv}, \code{H2_intradiario_benchmarks.csv} and \code{H2_curvas_intradiario.csv}.}
\end{table}

On 430 common sessions, intraday EWMA had QLIKE 0.0677 lower than close-to-close EWMA for the opening-to-close variance target. It reached 12.19\% annual volatility, closer to the 12\% target. Its higher exposure increased turnover. Net return was 7.59 percentage points lower, and annual CER was 10.77 percentage points lower. The CER ranking also held for $\gamma=1$ and $\gamma=5$.

\subsection{Matched exposure and drawdown}
The frozen signal trades when $|\widehat R^{OC}_{t+1}|>0.0035$. Ten trades emerged from 210 archived origins, seven with net gains. The 95\% Wilson interval for the hit rate is 39.7\% to 89.2\%.

\begin{table}[h!]
\centering\caption{Maximum drawdown and net return of the frozen signal.}\label{tab:drawdown}
\small
\begin{tabular*}{\textwidth}{@{\extracolsep{\fill}}lrrr@{}}
\toprule
Policy & Mean exposure & Max. drawdown & Net return \\
\midrule
Original signal & $4.78\%$ & $-1.77\%$ & $+2.94\%$ \\
Regime sizing & $2.05\%$ & $-0.77\%$ & $+1.16\%$ \\
Matched fixed sizing & $2.05\%$ & $-0.76\%$ & $+1.26\%$ \\
\bottomrule
\end{tabular*}
\fonte{\code{trading_sinal_congelado.csv}. Mean exposure includes days in cash.}
\end{table}

At matched average exposure, fixed sizing produced similar drawdown and higher net return than regime sizing. In 5,000 within-side permutations of the weights, the median return was $+1.24\%$. Regime sizing fell at the descriptive 2.6th percentile. Both controls are retrospective.

\subsection{Execution-cost sensitivity}
Keeping the ten trades and a fee of ten basis points per side, slippage of 0, 5, 10, and 20 basis points per side produced compounded net returns of $+3.97\%$, $+2.94\%$, $+1.91\%$, and $-0.10\%$. Solving $\prod_i[1+R_i^{\rm gross}-2(10+s)/10^4]=1$ gives break-even slippage of 19.49 basis points per side. This calculation holds the threshold and dates fixed. Separating fees, spread, and market impact would require observable quotes and order-size information.
\fonte{\code{auditoria_operacoes_exploratorio_2026.csv} and \code{auditoria_slippage_exploratorio_2026.csv}.}

\section{Discussion and conclusion}
The historical evidence for $H_1$ is inconclusive. For $H_2$, matched fixed exposure reproduces nearly all the drawdown reduction and retains more return. Execution-cost sensitivity explains why $H_3$ requires prospective prices. The ten trades provide an exploratory result with substantial uncertainty about the future hit rate.

\subsection{Empirical extensions}
With audited five-minute bars, realized variance would be
\[
 RV_t=\sum_{i=1}^{M_t}(\log P_{t,i}-\log P_{t,i-1})^2.
\]
HAR-RV, daily HAR, GARCH, and EWMA would then be evaluated against the same $RV_{t+1}$ targets. Both pilot sessions failed the coverage audit. For regime forecasts, a separate extension would compare extreme-event calibration and joint VaR--Expected Shortfall loss on new dates. A confirmatory HAR--EWMA comparison would fix the horizon, training windows, refit schedule, and primary test before observing new targets.

The signal remains frozen for prospective evaluation. Under the recorded ETF schedule, the reference entry is the 10:06 bar. The first review is scheduled for 30 directional signals or 500 origins. As of 7 November 2026, the cohort had zero observed trades.

\subsection{Reproducibility}
The private daily release manifest \code{V1_diaria.json} records 49 tests, 40 reproduced tables, and four reproduced artifacts. The public package includes the aggregate HAR--EWMA results and the script that generates them from a local forecast export.
\smallskip\noindent\textbf{References}\par
{\small
[1] Hamilton, J. D. (1989). A New Approach to the Economic Analysis of Nonstationary Time Series and the Business Cycle. \textit{Econometrica}, 57(2), 357--384. \href{https://doi.org/10.2307/1912559}{DOI}.\par
[2] Corsi, F. (2009). A Simple Approximate Long-Memory Model of Realized Volatility. \textit{Journal of Financial Econometrics}, 7(2), 174--196. \href{https://doi.org/10.1093/jjfinec/nbp001}{DOI}.\par
[3] Patton, A. J. (2011). Volatility forecast comparison using imperfect volatility proxies. \textit{Journal of Econometrics}, 160(1), 246--256. \href{https://public.econ.duke.edu/~ap172/Patton_vol_proxies_JoE_2011.pdf}{Paper}.\par
[4] Diebold, F. X.; Mariano, R. S. (1995). Comparing Predictive Accuracy. \textit{Journal of Business \& Economic Statistics}, 13(3), 253--263. \href{https://doi.org/10.1080/07350015.1995.10524599}{DOI}.\par
[5] DeMiguel, V., Garlappi, L., and Uppal, R. (2009). Optimal Versus Naive Diversification. \textit{Review of Financial Studies}, 22(5), 1915--1953. \href{https://papers.ssrn.com/sol3/papers.cfm?abstract_id=1376199}{Paper}.\par
}
\end{document}
